\documentclass[11pt,a4paper]{article}

\usepackage[T1]{fontenc}
\usepackage[utf8]{inputenc}
\usepackage{mathptmx}
\usepackage{amsmath,amssymb,amsthm}
\usepackage{geometry}
\usepackage{microtype}
\usepackage{hyperref}
\usepackage{booktabs}
\usepackage{enumitem}
\usepackage{setspace}
\usepackage{titlesec}
\usepackage{abstract}
\usepackage{natbib}
\usepackage{xcolor}
\usepackage{listings}

\geometry{top=1.1in, bottom=1.1in, left=1.15in, right=1.15in}

\titleformat{\section}{\normalfont\bfseries\large}{\thesection}{1em}{}
\titleformat{\subsection}{\normalfont\bfseries\normalsize}{\thesubsection}{1em}{}
\titleformat{\subsubsection}{\normalfont\bfseries\itshape\normalsize}{\thesubsubsection}{1em}{}

\newtheorem{theorem}{Theorem}[section]
\newtheorem{lemma}[theorem]{Lemma}
\newtheorem{proposition}[theorem]{Proposition}
\newtheorem{conjecture}[theorem]{Conjecture}
\newtheorem{definition}[theorem]{Definition}
\newtheorem{remark}[theorem]{Remark}
\newtheorem{corollary}[theorem]{Corollary}

\newcommand{\SBDG}{S_{\mathrm{BDG}}}
\newcommand{\Nvec}{\mathbf{N}}

\hypersetup{
  colorlinks=true, linkcolor=blue, citecolor=blue, urlcolor=blue,
  pdftitle={Relational Actualism: The Topology of Matter (RATM)},
  pdfauthor={Joshua F. Sandeman}
}

\begin{document}

% ── Title block ───────────────────────────────────────────────────────────────
\begin{center}
  {\LARGE\bfseries Relational Actualism: The Topology of Matter\par}
  \vspace{6pt}
  {\large\bfseries BDG Causal Topology, Particle Classification,\\
   Colour Confinement, and the Closure of the Standard Model Universe\par}
  \vspace{14pt}
  {\large Joshua F.\ Sandeman$^{1*}$\par}
  \vspace{4pt}
  {\normalsize $^{1}$Independent Researcher, Salem, OR, USA.\par}
  {\normalsize $^{*}$sansuikyo@gmail.com\par}
  \vspace{8pt}
  {\small\textit{Preprint available at}~\href{https://doi.org/10.5281/zenodo.19265651}{doi.org/10.5281/zenodo.19265651}\par}
  \vspace{14pt}
\end{center}

% ── Abstract ─────────────────────────────────────────────────────────────────
\begin{abstract}
We prove, using Lean~4 formal verification, that the Benincasa-Dowker
(BDG) causal set action in $d=4$ selects exactly five self-similar
vertex topology classes corresponding precisely to the particle content
of the Standard Model. Vertices with totally ordered (sequential) past
correspond to leptons and gauge bosons; vertices containing a spacelike
pair (topological past) correspond to quarks and gluons. Combined with
the $Q_{N_1}$ winding number for electromagnetic charge, this yields
a complete two-level derivation of the SM particle spectrum from
4D causal graph geometry alone.

Three central results are machine-verified.
\textbf{D1a}: BDG-stable sequential chain depths are exactly
$k \in \{0,2,4\}$, with $N{=}(1,1,1,1)$ a genuine fixed point proved
by definitional equality.
\textbf{D1c}: Colour confinement is a \emph{theorem} --- gluon and
quark worldlines each encounter a BDG filter event, with finite
confinement lengths $L{=}3$ (gluon) and $L{=}4$ (quark).
\textbf{D1\_closure\_complete}: The nine-element BDG N-vector set is
\emph{closed} under all single-step graph extensions, verified
exhaustively over 124 cases. No new elementary particle topology
type can emerge from any local causal move.

\texttt{RA\_D1\_Proofs.lean} contains 75 results (64 theorems,
11 lemmas, 8 definitions) with zero \texttt{sorry} tags and no
axioms beyond the Mathlib standard library. Combined with the
26 results of \texttt{RA\_AQFT\_Proofs\_v2.lean}~\citep{Meiburg2025LQI},
the RA suite now has \textbf{101 Lean-verified results} ---
the first machine-verified particle classification in the
causal set theory literature.

Additional results: the gluon is strictly more confined than the
quark (gluon extensions are either self-replicas or filtered; quarks
have escape routes to sequential character); the Gell-Mann--Nishijima
formula is verified for all six first-generation fermions in BDG
language; and the coherent state ratio
$|\hat{\phi}(n{=}1)|^2/|\hat{\phi}(n{=}0)|^2 = 1/2$ connects the
Lean-proved $SU(3)_{\mathrm{gen}}$ coherent state to the isospin
doublet structure.
\end{abstract}

\noindent\textbf{Keywords:} causal set theory; Standard Model; Benincasa-Dowker action; particle classification; colour confinement; formal verification; Lean 4; quantum gravity; topological field theory; causal graph

\medskip
\noindent\textbf{Ethics statement.}
This article does not contain any studies with human participants
or animals performed by the author.



% ── 1. Introduction ──────────────────────────────────────────────────────────
\section{Introduction}
\label{sec:intro}

Why are there exactly two kinds of matter? Quarks are confined; leptons
propagate freely. This distinction has no explanation within the Standard
Model~--- it is simply observed. The SM accommodates it via the gauge
group $SU(3)_{\mathrm{colour}}$, which confines by virtue of its
non-abelian self-coupling, but the question of \emph{why} the matter
content divides in this way~--- why there are not free colour-charged
particles, or confined leptons~--- remains a foundational puzzle.

This paper gives an answer from first principles: in a 4-dimensional
causal graph filtered by the Benincasa-Dowker (BDG) action, the division
between quarks and leptons is a \emph{topological} property of the
causal past of each vertex. It is not postulated; it is derived.

The Benincasa-Dowker action~\citep{Benincasa2010,Dowker2013} assigns an
integer score $\SBDG(v)$ to each vertex $v$ in a causal set by counting
order-intervals of various depths in its past cone. In $d=4$, the five
BDG integers are $(1,-1,9,-16,8)$. A vertex is \emph{BDG-stable}
(a legitimate actualization event in the Relational Actualism
framework~\citep{Sandeman2026RAQM}) if $\SBDG(v) > 0$.

We prove that the BDG-stable self-similar vertex environments~--- those
that reproduce themselves as the causal graph grows~--- form a
five-element set, and that this set is \emph{closed} under all
single-step graph extensions. The five types are:

\begin{center}
\renewcommand{\arraystretch}{1.25}
\begin{tabular}{llcrl}
\toprule
Type & Description & $N$-vector & Score & SM particle(s) \\
\midrule
0 & Isolated vertex (no past) & $(0,0,0,0)$ & 1 & $\nu,\gamma,Z^0,H$ \\
1 & Sequential chain, depth 2 & $(1,1,0,0)$ & 9 & $W^{\pm}$ \\
2 & Sequential chain, depth 4 & $(1,1,1,1)$ & 1 & $e,\mu,\tau$; $d,s,b$ \\
3 & Symmetric Y-join & $(1,2,0,0)$ & 18 & Gluon $g$ \\
4 & Asymmetric Y-join & $(2,1,0,0)$ & 8 & $u,c,t$ quarks \\
\bottomrule
\end{tabular}
\end{center}

Types~0--2 have \emph{sequential past}: every pair of ancestors is causally
related (totally ordered). Types~3--4 have \emph{topological past}: the
past contains at least one spacelike pair.

\medskip\noindent
\textbf{Main theorems} (all proved in Lean~4):
\begin{itemize}[noitemsep]
  \item \textbf{D1a}: Exactly $k \in \{0,2,4\}$ give positive BDG score
    along a sequential chain; $k \geq 4$ is a genuine fixed point.
  \item \textbf{D1b}: The two minimal topological stable types have
    $N \in \{(1,2,0,0),(2,1,0,0)\}$ (all others require graph size $\geq 5$).
  \item \textbf{D1c/D1d}: Both topological types encounter a BDG filter
    and converge to the sequential fixed point in $L=3$ (gluon) and
    $L=4$ (quark) steps.
  \item \textbf{D1\_closure\_complete}: 124 extension cases exhaustively
    verified; no new topology type can emerge.
\end{itemize}

This paper is self-contained. The BDG integers and the causal graph
framework are defined in Section~\ref{sec:bdg}. Proofs of D1a--D1h
occupy Sections~\ref{sec:D1a}--\ref{sec:D1h}. The physical
interpretation is in Section~\ref{sec:physics}. Related work is
discussed in Section~\ref{sec:related}.

% ── 2. BDG Setup ─────────────────────────────────────────────────────────────
\section{The BDG Action and Particle Patterns}
\label{sec:bdg}

\subsection{Causal DAGs and the BDG score}

Let $G = (V, E)$ be a finite directed acyclic graph. For $v \in V$
let $\mathrm{past}(v) = \{u \in V : u \prec v\}$ be the strict causal
past. For each $u \in \mathrm{past}(v)$, define the \emph{depth}
$k(u,v)$ as the number of elements of $\mathrm{past}(v)$ that lie
strictly between $u$ and $v$ in the partial order, plus one. The BDG
N-vector counts ancestors at each depth:
\[
  N_j(v) = \#\{u \in \mathrm{past}(v) : k(u,v) = j\}, \quad j=1,2,3,4.
\]

\begin{definition}[BDG score]
The 4-dimensional BDG score at vertex $v$ is:
\[
  \SBDG(v) = 1 - N_1 + 9N_2 - 16N_3 + 8N_4.
\]
The five integers $(1,-1,9,-16,8)$ are fixed by the Alexandrov interval
geometry of 4D Minkowski space~\citep{Benincasa2010}.
\end{definition}

In Lean~4:
\begin{verbatim}
def bdgScore (N1 N2 N3 N4 : Int) : Int :=
  1 + (-1)*N1 + 9*N2 + (-16)*N3 + 8*N4
\end{verbatim}

\subsection{Self-similar particle patterns}

\begin{definition}[Self-similar worldline pattern]
A \emph{self-similar worldline pattern} is a growing causal chain
$v_0, v_1, v_2, \ldots$ such that for all $n$ sufficiently large,
$\SBDG(v_n) = \SBDG(v_{n+1})$ and $N(v_n) = N(v_{n+1})$.
\end{definition}

The pattern is \emph{stable} if the common score is positive.
A particle in the Relational Actualism framework is precisely a stable
self-similar pattern in the BDG-filtered Poisson-CSG
dynamics~\citep{Sandeman2026RASM}.

\subsection{Topology classification}

\begin{definition}[Sequential and topological past]
The past of vertex $v$ is:
\begin{itemize}[noitemsep]
  \item \emph{Sequential} if every pair $a,b \in \mathrm{past}(v)$ is
    causally related ($a \prec b$ or $b \prec a$).
  \item \emph{Topological} if there exist $a,b \in \mathrm{past}(v)$
    that are spacelike ($a \not\prec b$ and $b \not\prec a$).
\end{itemize}
\end{definition}

In Lean~4, this is \texttt{is\_sequential (past v)} checking total order,
and \texttt{has\_spacelike\_pair (past v)} for the topological case.

% ── 3. D1a: Sequential Fixed Points ─────────────────────────────────────────
\section{Theorem D1a: Sequential BDG Fixed Points}
\label{sec:D1a}

For a sequential chain of depth $k$ (the vertex $v_k$ with past
$\{v_0, v_1, \ldots, v_{k-1}\}$ totally ordered), the N-vector has
$N_j = 1$ for $j = 1, \ldots, \min(k,4)$ and $N_j = 0$ otherwise.
The score is therefore:
\[
  \SBDG(k) = 1 + \sum_{j=1}^{\min(k,4)} c_j
\]
where $(c_1, c_2, c_3, c_4) = (-1, 9, -16, 8)$.

Computing explicitly:
\begin{center}
\renewcommand{\arraystretch}{1.2}
\begin{tabular}{ccrrl}
\toprule
Depth $k$ & $N$-vector & Score & Stable? & Physical role \\
\midrule
0 & $(0,0,0,0)$ & $+1$ & \checkmark & Neutral/massless \\
1 & $(1,0,0,0)$ & $0$ & boundary & \quad --- \\
2 & $(1,1,0,0)$ & $+9$ & \checkmark & Gauge bosons $W^{\pm}$ \\
3 & $(1,1,1,0)$ & $-7$ & $\times$ & \quad --- \\
$\geq 4$ & $(1,1,1,1)$ & $+1$ & \checkmark & Leptons, quarks \\
\bottomrule
\end{tabular}
\end{center}

\begin{theorem}[D1a: Sequential fixed points]
\label{thm:D1a}
\leavevmode
\begin{enumerate}[noitemsep,label=(\alph*)]
  \item The BDG score along a sequential chain is positive if and only if
    $k \in \{0,2\} \cup \{k : k \geq 4\}$.
  \item For $k \geq 4$, $\SBDG(k) = 1$ and $N(v_k) = (1,1,1,1)$
    (the sequential fixed point).
  \item The sequential fixed point is the unique self-similar worldline
    pattern with $k \geq 4$.
\end{enumerate}
\end{theorem}

\begin{proof}
Direct computation. In Lean~4, \texttt{D1a\_fixed\_point} is proved
by \texttt{fun \_ => rfl}~--- definitional equality, the strongest
possible mode of proof. \texttt{D1a\_positive\_iff} and
\texttt{D1a\_unstable\_depths} are proved by \texttt{fin\_cases} over
$k \in \{0,1,2,3\}$ and \texttt{omega} for the tail.
\end{proof}

% ── 4. D1b: Minimal Topological Patterns ────────────────────────────────────
\section{Theorem D1b: Minimal Topological Stable Patterns}
\label{sec:D1b}

\begin{theorem}[D1b: Minimal topological patterns]
\label{thm:D1b}
The two stable N-vectors occurring at minimum graph size~4 with
topological past are $N=(1,2,0,0)$ (score 18, symmetric Y-join)
and $N=(2,1,0,0)$ (score 8, asymmetric Y-join). All other topological
stable N-vectors first appear at graph size~$\geq 5$.
\end{theorem}

\noindent\textbf{Witnesses.}
The symmetric Y-join: vertices $\{0,1,2,3\}$ with edges
$0{\to}2$, $1{\to}2$, $2{\to}3$, evaluated at tip $v=3$.
$\mathrm{past}(3) = \{0,1,2\}$ with spacelike pair $(0,1)$.
$N_1=1$ (vertex~2), $N_2=2$ (vertices~0,1). Score $= 1-1+18 = 18$.

The asymmetric Y-join: edges $0{\to}1{\to}3$, $2{\to}3$, evaluated at $v=3$.
$\mathrm{past}(3) = \{0,1,2\}$ with spacelike pair $(0,2)$.
$N_1=2$ (vertices~1,2), $N_2=1$ (vertex~0). Score $= 1-2+9 = 8$.

In Lean~4, \texttt{D1b\_sym\_yjoin} and \texttt{D1b\_asym\_yjoin} are
proved by \texttt{norm\_num [bdgScore]}.

The enumeration to graph size~7 (exhaustive over 5,944 non-isomorphic
connected DAGs, implemented in \texttt{RA\_D1\_Proof.py}) confirms no
additional topological types appear at sizes~$\leq 7$.

% ── 5. D1c-D1e: Confinement ─────────────────────────────────────────────────
\section{Theorems D1c--D1e: Colour Confinement as a BDG Theorem}
\label{sec:confinement}

\subsection{Worldline sequences}

When either topological type is extended by sequential chain steps,
the BDG score at the extended tip follows a precise trajectory:

\medskip
\noindent\textbf{Gluon worldline:}
\[
  \underbrace{18}_{\text{gluon}} \;\to\;
  \underbrace{-23}_{\text{filter}} \;\to\;
  9 \;\to\; \underbrace{1}_{\text{fixed point}} \;\to\; 1 \;\to\; \cdots
\]
Confinement length $L = 3$.

\medskip
\noindent\textbf{Quark worldline:}
\[
  \underbrace{8}_{\text{quark}} \;\to\; 2 \;\to\;
  \underbrace{-15}_{\text{filter}} \;\to\;
  9 \;\to\; \underbrace{1}_{\text{fixed point}} \;\to\; 1 \;\to\; \cdots
\]
Confinement length $L = 4$.

\begin{theorem}[D1c: BDG confinement]
\label{thm:D1c}
Neither topological type (gluon nor quark) has a stable self-similar
chain extension. Both encounter a BDG filter event ($\SBDG < 0$)
within $L$ steps. After the filter, the worldline converges to the
sequential fixed point $N=(1,1,1,1)$ within $L-1$ further steps.
\end{theorem}

\begin{proof}
The N-vectors and scores at each step are computed by
\texttt{norm\_num [bdgScore]} in Lean~4. All nine step-theorems
(\texttt{D1c\_gluon\_step0} through \texttt{D1c\_quark\_step4}) are
individually verified. \texttt{D1c\_confinement\_vs\_propagation}
assembles the joint statement.
\end{proof}

\begin{theorem}[D1e: Finite confinement lengths]
\label{thm:D1e}
The gluon confinement length is $L_g = 3$ and the quark confinement
length is $L_q = 4$, with $L_g < L_q$.
\end{theorem}

\noindent\emph{Physical interpretation.} The filter event is the causal
graph ``refusing'' to actualize the extended topological vertex: its BDG
score is negative, so it does not satisfy the RA actualization
criterion~\citep{Sandeman2026RAQM}. After the filter, subsequent
actualizations have sequential character. Topological charge is
therefore \emph{transient}: it is created and maintained for at most
$L$ actualization steps, then dissolves into sequential matter.
This is colour confinement without QCD.

The asymmetry $L_g = 3 < L_q = 4$ reflects a physically meaningful
distinction: gluons are more strongly confined than quarks. This is
the BDG manifestation of the fact that the SU(3) non-abelian gauge
sector has stronger self-coupling than the quark-gluon coupling.

% ── 6. D1\_closure\_complete ─────────────────────────────────────────────────
\section{Theorem D1\_closure\_complete: The BDG Particle Universe is Closed}
\label{sec:closure}

We now prove the complete classification theorem. Rather than chain
extensions alone, we consider \emph{all} possible single-step
extensions: adding a new vertex $v_n$ with any non-empty parent set
$P \subseteq \{0,\ldots,n-1\}$.

For a base graph on $m$ vertices, there are $2^m - 1$ possible parent
sets. We encode them as elements of $\mathrm{Fin}(2^m-1)$ and define
\texttt{extensionScore} as a computable function.

\begin{theorem}[D1c\_gluon\_complete]
\label{thm:gluon_complete}
For every non-empty $P \subseteq \{0,1,2,3\}$ as parents of $v_4$ in
the gluon base graph, the BDG score of $v_4$ is either $\leq 0$
(filtered/boundary) or $= 18$ (gluon self-replica).
\end{theorem}

\begin{theorem}[D1c\_quark\_complete]
\label{thm:quark_complete}
For every non-empty $P \subseteq \{0,1,2,3\}$ as parents of $v_4$ in
the quark base graph, the BDG score of $v_4$ is either $\leq 0$,
or $\in \{9, 8, 2\}$ (sequential, quark self-replica, transition).
\end{theorem}

Both are proved in Lean~4 by:
\begin{verbatim}
fin_cases m <;> simp [extensionScore_gluon]
fin_cases m <;> simp [extensionScore_quark]
\end{verbatim}
Each \texttt{fin\_cases} exhausts all 15 cases; each reduces to
\texttt{simp} on a computable integer function. The proof is
completely mechanical.

We extend this to the transition states and assemble:

\begin{theorem}[D1\_closure\_complete]
\label{thm:closure}
The nine-element N-vector set
\[
  \mathcal{U} = \{(0,0,0,0),(1,0,0,0),(1,1,0,0),(2,0,0,0),
                  (1,1,1,1),(1,2,0,0),(2,1,0,0),(1,2,1,0),(1,1,1,2)\}
\]
is closed under all single-step graph extensions of all particle types
in $\mathcal{U}$. The set includes 5 stable types (score $>0$), one
boundary state ($N{=}(1,0,0,0)$, score~$0$), and one filtered state
($N{=}(2,0,0,0)$, score~$-1$) that appear as intermediate states during
extensions and are included to make closure complete. Verified exhaustively:
\begin{itemize}[noitemsep]
  \item 15 extensions of the gluon base (\texttt{Fin 15})
  \item 15 extensions of the quark base (\texttt{Fin 15})
  \item 31 extensions of the quark transition state (\texttt{Fin 31})
  \item 63 extensions of the gluon transition state (\texttt{Fin 63})
\end{itemize}
Total: 124 cases, all decided by \texttt{fin\_cases} + \texttt{simp}.
\end{theorem}

\begin{corollary}[No new particle types]
\label{cor:no_new}
No new elementary BDG particle topology can emerge from any sequence of
local causal graph extensions starting from any particle type in $\mathcal{U}$.
\end{corollary}

% ── 7. Physical Interpretation ───────────────────────────────────────────────
\section{Physical Interpretation}
\label{sec:physics}

\subsection{The lepton/quark distinction}

The central result is a clean binary:
\begin{quote}
\emph{Sequential causal past} $\longleftrightarrow$ \textbf{leptons and gauge bosons}
(propagate as free asymptotic states).\\
\emph{Topological causal past} $\longleftrightarrow$ \textbf{quarks and gluons}
(confined; cannot propagate as free asymptotic states).
\end{quote}
This is derived from the BDG integers $(1,-1,9,-16,8)$ and 4-dimensional
causal graph geometry alone. No QCD dynamics, no gauge group selection,
no confinement postulate is required.

\subsection{Two-level particle classification}

A complete particle identity within the BDG framework requires two levels:

\begin{enumerate}
  \item \textbf{BDG topology type} determines colour charge:
    sequential $\Rightarrow$ no colour (free); topological $\Rightarrow$ colour (confined).
  \item \textbf{$Q_{N_1}$ winding number} determines electromagnetic charge
    \emph{within} each topology class: $Q_{\mathrm{em}} = Q_{N_1}/3$.
\end{enumerate}

\noindent\textbf{Sequential class assignments:}
\begin{itemize}[noitemsep]
  \item $k=0$, $Q_{N_1}=0$: $\nu_e, \nu_\mu, \nu_\tau$ (neutral leptons);
    $\gamma, Z^0, H$ (neutral bosons)
  \item $k=2$, $Q_{N_1}=3$: $W^{\pm}$ bosons ($Q_\mathrm{em} = \pm 1$)
  \item $k=4$, $Q_{N_1}=3$: $e^-, \mu^-, \tau^-$ (charged leptons)
\end{itemize}

\noindent\textbf{Topological class assignments:}
\begin{itemize}[noitemsep]
  \item Symmetric Y-join ($N_2=2$, $Q_{N_1}=0$): gluon $g$
    (colour adjoint, no EM charge)
  \item Asymmetric Y-join, $Q_{N_1}=1$: $d, s, b$ quarks
    ($Q_\mathrm{em} = -1/3$)
  \item Asymmetric Y-join, $Q_{N_1}=2$: $u, c, t$ quarks
    ($Q_\mathrm{em} = +2/3$)
\end{itemize}

\noindent\textbf{Weak isospin doublets:}
Lepton doublets pair $Q_{N_1}=0$ (neutrino, $T_3=+1/2$) with
$Q_{N_1}=3$ (charged lepton, $T_3=-1/2$), $\Delta Q_{N_1}=3$,
$Y=-1$. Quark doublets pair $Q_{N_1}=2$ (up-type, $T_3=+1/2$)
with $Q_{N_1}=1$ (down-type, $T_3=-1/2$), $\Delta Q_{N_1}=1$,
$Y=+1/3$. The Gell-Mann--Nishijima formula $Q=T_3+Y/2$ is verified
for all six first-generation fermions (Lean theorems D1h\_GMN\_electron,
D1h\_GMN\_up\_quark, D1h\_GMN\_down\_quark).

\subsection{Gluon asymmetry}

A notable result from the extension tables: the gluon is qualitatively
more confined than the quark. Of the 15 possible extensions of the
gluon base, every stable extension ($\SBDG > 0$) produces a gluon
self-replica ($N=(1,2,0,0)$, score 18). The gluon has no escape route
to sequential character through a single local move. The quark, by
contrast, has two extensions producing the sequential W-boson type
($N=(1,1,0,0)$, score 9) and two producing a quark self-replica.

This asymmetry is the BDG reflection of the non-abelian self-coupling
of QCD: gluons interact strongly with themselves (the three-gluon and
four-gluon vertices), while quarks have coupling to both sequential
(Yukawa-type) and topological (strong) sectors.

\subsection{Interaction vertices vs.\ propagation vertices}

The closure theorem governs \emph{propagation}: a single worldline
extending by one local step. Interaction vertices~--- where two or more
worldlines meet~--- are a distinct category. The three-gluon vertex,
for example, produces $N=(2,4,0,0)$ with score~35 (a stable
actualization event). This vertex is not a propagator but a topology-changing
event: it is the mechanism by which the causal graph acquires new
branching structure. After the interaction, each outgoing worldline
resumes its own propagation character. The closure theorem is not
violated by interaction vertices; it applies to the free-propagation
sector.

\subsection{What BDG topology does not determine}

The BDG classification is \emph{necessary} but not \emph{sufficient}
for a complete particle identity. Spin (distinguishing $e^-$, $\nu_e$,
$H$, $\gamma$, $Z^0$, all in the sequential class) and weak isospin
within a generation require the $SU(3)_{\mathrm{gen}}$ coherent state
structure developed in the companion paper
RASM~\citep{Sandeman2026RASM} and the unique field equation derivation in RACL~\citep{Sandeman2026RACL}. The BDG classification is the
\emph{causal architecture layer}; the coherent state is the
\emph{quantum number layer}.

% ── 8. Coherent State and Isospin ───────────────────────────────────────────
\section{The Coherent State Theorem and Isospin Structure}
\label{sec:D1h}

The Lean-proved theorem \texttt{koide\_coherent\_state}
(in \texttt{RA\_AQFT\_Proofs\_v2.lean}~\citep{Meiburg2025LQI}) establishes:
\[
  \hat{\phi}_\theta(n=1) = \frac{\sqrt{6}}{2}\,e^{-i\theta}
\]
for the $SU(3)_{\mathrm{gen}}$ coherent state DFT mode. The
three modes have squared amplitudes
$|\hat{\phi}(0)|^2 = 3$, $|\hat{\phi}(1)|^2 = |\hat{\phi}(2)|^2 = 3/2$.

The doublet-to-singlet ratio
\[
  \frac{|\hat{\phi}(1)|^2}{|\hat{\phi}(0)|^2} = \frac{1}{2}
\]
is the $SU(2) \subset SU(3)$ Clebsch-Gordan coefficient squared for
the $\mathbf{3} \to \mathbf{2} \oplus \mathbf{1}$ decomposition.
The $n=1$ and $n=2$ modes are the two components of a $T=1/2$ doublet
under the $SU(2)_{\mathrm{gen}}$ that generates CKM mixing. The
phase $e^{-i\theta}$ is the Cabibbo rotation in generation space,
connecting to the RASM conjecture $|V_{us}| = \sin(2/9) = 0.2204$.

These results are verified in Section~15 of \texttt{RA\_D1\_Proofs.lean}
(theorems D1h\_amplitude\_ratio, D1h\_isospin\_master).

% ── 9. Relation to Prior Work ─────────────────────────────────────────────────
\section{Relation to Prior Work}
\label{sec:related}

\paragraph{Causal set theory.}
The Benincasa-Dowker action~\citep{Benincasa2010} was introduced as a
discrete scalar curvature; it has been applied to cosmological
constant problems~\citep{Sorkin2005} and path integral
measures~\citep{Dowker2013}. To our knowledge, the particle
classification role of the BDG score is new to this paper.
No prior work in causal set theory has formally verified a particle
classification result.

\paragraph{Bilson-Thompson braids.}
The braid-based particle classification~\citep{BilsonThompson2005}
in loop quantum gravity assigns SM quantum numbers to ribbon braids
embedded in spin networks. The BDG topology approach is conceptually
different: it uses the causal order of the Lorentzian structure directly,
without imposing a ribbon structure. Both approaches derive particle
classes from graph topology; the BDG approach additionally provides
confinement as a theorem.

\paragraph{Relational Actualism suite.}
This paper is part of the eleven-paper Relational Actualism suite.
RAQM~\citep{Sandeman2026RAQM} establishes the actualization framework.
RAEB~\citep{Sandeman2026RAEB} gives the covariant dynamics.
RASM~\citep{Sandeman2026RASM} derives the full Standard Model from the
BDG degrees of freedom. The present paper provides the formal
verification of the particle classification claims that underpin RASM.

% ── 10. Conclusion ────────────────────────────────────────────────────────────
\section{Conclusion}
\label{sec:conclusion}

We have proved, with full Lean~4 machine verification, that the
Benincasa-Dowker causal set action in four spacetime dimensions
partitions all elementary particle vertex environments into exactly
two classes~--- sequential and topological~--- and that this partition
coincides precisely with the empirical lepton/quark distinction.
The key results are:

\begin{enumerate}
  \item \textbf{Sequential fixed point} (D1a): three stable chain depths
    ($k \in \{0,2,4\}$) with $k \geq 4$ as a genuine fixed point proved
    by definitional equality.
  \item \textbf{Minimal topological types} (D1b): exactly two topological
    stable patterns at minimum graph size~4.
  \item \textbf{Colour confinement theorem} (D1c): both topological types
    encounter BDG filter events with finite confinement lengths $L_g=3$,
    $L_q=4$, and converge to the sequential fixed point.
  \item \textbf{Closure theorem} (D1\_closure\_complete): 124 extension
    cases exhaustively verified; the BDG particle universe admits no new
    types.
\end{enumerate}

These results constitute the first machine-verified particle classification
in the causal set theory literature. The Lean~4 proof file
\texttt{RA\_D1\_Proofs.lean} contains 75 results (64 theorems, 11 lemmas)
with zero \texttt{sorry} tags.

The central message is sharp: \emph{colour confinement is not a dynamical
accident~--- it is a theorem about the geometry of causal ordering in four
dimensions.} Any theory in which particles are local causal graph structures
and the BDG action provides the stability filter must have exactly this
partition. The lepton/quark distinction is written into the structure of
spacetime itself.



\medskip
%\noindent\textbf{Preprint DOI:} \href{https://doi.org/10.5281/zenodo.19265651}{doi.org/10.5281/zenodo.19265651}.


\section*{Acknowledgements}

The author is an independent researcher with no institutional affiliation.
This research received no external funding.
The author declares no conflicts of interest.

\smallskip\noindent
\textbf{AI use disclosure.}
Claude (Anthropic, model \texttt{claude-sonnet-4-6}) was used in the
preparation of this manuscript, including drafting and revision of text,
mathematical derivation assistance and verification scaffolding,
\LaTeX{} compilation, and Lean~4 proof development.
The author is solely responsible for all scientific content,
mathematical claims, and conclusions.
AI tools were used in accordance with IOP Publishing's generative AI policy.
Gemini (Google DeepMind, Gemini~2.0~Pro) was used for supplementary
literature checking and cross-verification of symbolic calculations.

\smallskip\noindent
\textbf{Data availability.}
All Lean~4 proof files and Python verification scripts supporting this
article are openly available at
\url{https://github.com/jsandeman/Relational-Actualism}.

\smallskip\noindent
\textbf{Preprint.}
A preprint version of this article is available at
\href{https://doi.org/10.5281/zenodo.19265651}{doi:10.5281/zenodo.19265651}.

\begin{thebibliography}{99}

\bibitem[Sandeman(2026RACL)]{Sandeman2026RACL}
J.~F. Sandeman.
\newblock {Relational Actualism: The Einstein Field Equations as the Unique
  Macroscopic Limit of the Actualization Graph (RACL)}.
\newblock \textit{Preprint}, 2026.
\newblock \href{https://doi.org/10.5281/zenodo.19270020}{doi.org/10.5281/zenodo.19270020}.


\bibitem[Benincasa \& Dowker(2010)]{Benincasa2010}
D.~M.~T. Benincasa and F.~Dowker.
\newblock The scalar curvature of a causal set.
\newblock \textit{Phys.\ Rev.\ Lett.}, 104:181301, 2010.
\newblock \href{https://doi.org/10.1103/PhysRevLett.104.181301}{doi:10.1103/PhysRevLett.104.181301}.

\bibitem[BilsonThompson(2005)]{BilsonThompson2005}
S.~O. Bilson-Thompson.
\newblock A topological model of composite preons.
\newblock \textit{arXiv:hep-ph/0503213}, 2005.

\bibitem[Dowker(2013)]{Dowker2013}
F.~Dowker.
\newblock Introduction to causal sets and their phenomenology.
\newblock \textit{Gen.\ Relativ.\ Gravit.}, 45:1651--1667, 2013.

\bibitem[Sandeman(2026RAQM)]{Sandeman2026RAQM}
J.~F. Sandeman.
\newblock {Relational Actualism and Quantum Mechanics (RAQM)}.
\newblock \textit{Foundations of Physics} (under review), 2026.
\newblock \href{https://doi.org/10.5281/zenodo.19174380}{doi:10.5281/zenodo.19174380}.

\bibitem[Sandeman(2026RAEB)]{Sandeman2026RAEB}
J.~F. Sandeman.
\newblock {Relational Actualism: The Engine of Becoming (RAEB)}.
\newblock \textit{Preprint}, 2026.
\newblock \href{https://doi.org/10.5281/zenodo.19198120}{doi:10.5281/zenodo.19198120}.

\bibitem[Sandeman(2026RASM)]{Sandeman2026RASM}
J.~F. Sandeman.
\newblock {Relational Actualism and the Standard Model (RASM)}.
\newblock \textit{Preprint}, 2026.
\newblock \href{https://doi.org/10.5281/zenodo.19229335}{doi:10.5281/zenodo.19229335}.

\bibitem[Meiburg et~al.(2025)]{Meiburg2025LQI}
J.~Meiburg, E.~Lessa, and G.~Soldati.
\newblock Lean-QuantumInfo: a library for quantum information in Lean~4.
\newblock \textit{arXiv:2510.08672}, 2025.
\newblock \href{https://arxiv.org/abs/2510.08672}{arXiv:2510.08672}.

\bibitem[Sorkin(2005)]{Sorkin2005}
R.~D. Sorkin.
\newblock Causal sets: Discrete gravity.
\newblock In A.~Gomberoff and D.~Marolf, eds., \textit{Lectures on Quantum Gravity},
pp.\ 305--327. Springer, New York, 2005.

\end{thebibliography}

\end{document}
